\SetPicSubDir{ch3-method}
\label{sec:Method}
This study explores the occupant comfort performance of \ac{occ} employing three types of \ac{gcm} based on different occupancy resolution and the effect of occupancy dynamics on the control.
To examine these purposes, agentic-simulation framework is used based on field data collected in a real-world office building.
The overall procedural flow of the study is illustrated in \autoref{fig:ex_flow}.
\begin{figure}[pos=htbp!]
    \centering
    \includegraphics[width=0.97\linewidth]{\Pic{png}{SimulationFlow_small}}
    \caption{Overall procedure flow of the study}
    \label{fig:ex_flow}
\end{figure}
% \begin{figure}[pos=htbp!]
%     \centering
%     \includegraphics[width=0.97\linewidth]{\Pic{png}{procedure}}
%     \caption{Overall procedure flow of the study}
%     \label{fig:ex_flow}
% \end{figure}

\subsection{Data Collection}
\label{subsec:DataCollect}
For the simulation, three primary types of field data were collected in the target building: real-time occupancy logs, subjective \ac{tcs} and environmental measurement logs.
For \ac{tcs} and indoor environment measurement, data collected and described in~\cite{horikoshi_double_2026} are used.
This subsection therefore describes the analysis target building and the procedure for location data tracking.

\subsubsection{Target Building, Participant Recruitment and Ethical Considerations}
\label{sec:TargetBuil}
The data collection was conducted at The GEAR, an R\&D and administrative office building located in Singapore and operated by Kajima Corporation, from 29 September to 24 October 2025 (4 weeks).
The 3rd to 6th floors, which accommodate a mix of shared-office and static seating work styles depending on departmental roles, were the target areas for data collection.
The target zones are served by a central air-conditioning system consisting of air-handling units and dedicated outdoor air-handling units, using variable air volume units capable of individual zone temperature and ventilation control.
\qty{45} office workers who regularly work in the designated target rooms were invited to participate.
Considering the completeness of the collected survey answers and location tracking records, data from \qty{39} participants were used in the analysis.
This study has been approved by the National University of Singapore Institutional Review Board (NUS-IRB Reference Code: NUS-IRB-2025-498).
Informed consent was obtained from all participants prior to the study.
All location and survey data were anonymized by the building administrator before being accessed by the research team to protect privacy. 

% \begin{table*}[pos=htbp!]
%     \centering
%     \small
%     % \TableSetup{4.0pt}{1.08}
%     \caption{Target rooms used for field data collection.}
%     \begin{tabular}{@{}p{1.45cm}p{4.00cm}M{0.80cm} N{2.0cm} M{2.2cm} M{1.70cm}@{}}
    \toprule
    \TableHeader{Room} & \TableHeader{Occupants' job function} & \TableHeader{Level} & \TableHeaderBreak{Floor area\TableHeaderLineBreak \((\mathrm{m^2})\)} & \TableHeaderBreak{Seating\TableHeaderLineBreak capacity} & \TableHeaderBreak{Survey\TableHeaderLineBreak participants} \\
    \midrule
    Office 3A & Facility management & 3 & 70.34 & 12 & 2 \\
    Office 3B & Business development & 3 & 43.08 & 6 & 3 \\
    Office 5A & Real estate development & 5 & 221.40 & 14 & 8 \\
    Office 5B & Building construction & 5 & 213.73 & 20 & 10 \\
    Office 5C & Real estate development & 5 & 221.40 & 15 & 9 \\
    Office 5D & Research \& development & 5 & 163.56 & 18 & 10 \\
    Office 6A & Interior design & 6 & 68.85 & 5 & 3 \\
    \bottomrule
\end{tabular}

%     \label{tab:target_rooms}
% \end{table*}

% \subsubsection{Thermal Comfort Survey}
% Subjective \ac{tcv} survey answers were collected via a web survey during work hours. 
% Items and scales follow ANSI/ASHRAE Standard 55:2023, Appendix~L \cite{ashrae_ansiashrae_2022}.
% The survey contents are summarized in \autoref{table:comfsurvey}.
% Responses within 20 minutes of room entry were excluded to avoid thermal transient effects.
% Participants were prompted up to six times per workday: three times in the morning and three times in the afternoon.
% Over the one-month observation period, \qty{2608} survey responses were collected in total.

% \subsubsection{Indoor Environmental Measurements}
% Indoor environmental data was monitored in parallel with \ac{tcv} survey.
% The primary measurements were \ac{ti} and \ac{rh}, obtained from compact sensors (HIOKI E.E. Corporation, Nagano, Japan) installed in the target rooms.
% For rooms with two to three compact sensors, their mean value was used as a representative room-level value to reduce sensor-specific measurement bias.
% To ensure sufficient variability of survey answers, \ac{ti} setpoints were cycled between \qtyrange{22}{27}{\celsius} in \qty{1}{\celsius} steps, with a randomized order across days.  
% On selected days, morning and afternoon setpoints differed to expand diversity.
% In addition, measured air speed (\acs{v}) and mean radiant temperature (\acs{tr}) were recorded periodically as diagnostic checks rather than as main focus of the analysis.

\subsubsection{Individual-level location tracking log}
Anonymized individual-level location records were obtained from a face-recognition system installed at room entrances and in corridors.
The timestamped zone-level records were reconstructed into stay logs, where each entry represents a continuous period of presence by one participant in one room and contains an anonymized participant ID, room ID, start time, and end time.
These stay logs provide the basis for constructing location-log agents in the simulation.
The detailed procedure for this reconstruction is described in Appendix~\ref{app:occupancy_reconstruction}.

\subsection{Occupancy-log-based agentic comfort simulation}
\subsubsection{Simulation Design}
Using the collected \ac{tcs} responses, environmental measurements, and reconstructed stay logs, two agentic components were prepared for each participant: a \ac{pcm} and a location-log agent.
The \ac{pcm} represents the participant's thermal preference response to indoor temperature, whereas the location-log agent preserves the participant's observed room-level presence pattern over the study period.
In the simulation, these two components were intentionally separated so that comfort-response characteristics could be randomly reassigned to observed occupancy patterns.
A Monte Carlo-style resampling simulation was then conducted by combining sampled location-log agents with randomly assigned \acp{pcm}.
% At each control timestep, the location-log agents determined the simulated room occupant composition.
% Based on the simulated occupant compositions, three \ac{gcm}-based \ac{occ} policies were evaluated from comfort and control perspectives.

For agentic simulations, 
\begin{table*}[pos=htbp!]
    \centering
    \small
    % \TableSetup{4.0pt}{1.08}
    \caption{Target rooms used for field data collection and simulation cases}
    \begin{tabular}{@{}p{2.6cm}p{3.8cm}>{\centering\arraybackslash}m{2.2cm}>{\centering\arraybackslash}m{1.8cm}>{\centering\arraybackslash}m{1.8cm}@{}}
    \toprule
    \textbf{Room} &
    \makecell[c]{\textbf{Occupants'}\\\textbf{job function}} &
    \makecell[c]{\textbf{Seating type}} &
    \makecell[c]{\textbf{Floor area}\\\textbf{\((\mathrm{m^2})\)}} &
    \makecell[c]{\textbf{Seating}\\\textbf{capacity}} \\
    \midrule
    Office 5A & Real estate development & Assigned & 221.40 & 15 \\
    Office 5B & Research \& development & Assigned & 163.56 & 18 \\
    Shared Office 5A & Real estate development & Shared & 221.40 & 14 \\
    Shared Office 5B & Building construction   & Shared & 213.73 & 20 \\
    \bottomrule
\end{tabular}

    \label{tab:SimuTarget}
\end{table*}

\subsubsection{Personal Comfort Model}
Individual \acp{pcm} are developed for each participant to predict thermal preference based on room-level \ac{ti}.
This input parameter selection reflects the intention to represent the actual controllable \ac{hvac} system parameters in the target buildings used for data collection.
The \ac{pcm}s are developed using ordered multinomial logistic classification, utilized in similar agentic-based thermal comfort simulations~\cite{jung_comparative_2019}\cite{jung_energy_2020}, motivated by its interpretable representation of thermal comfort probabilities across input parameter values.
The model predicts the probability distribution over three thermal preference classes: $S_{th} \in \{+1 (\text{Warmer}), 0 (\text{No change}), -1 (\text{Cooler})\}$, using \ac{ti} as the primary feature.
The temperature-dependent ``No change'' probability curve of an individual \ac{pcm} $i$ is defined as
\begin{equation}
p_i(T)
\equiv
P_i\!\left(
\underbrace{S_{th}=0}_{\text{No change}}
\,\middle|\, T
\right).
\label{eq:pcm_probability}
\end{equation}
Here, $T$ denotes indoor air temperature and $i$ identifies the individual \ac{pcm}.

\subsubsection{Monte Carlo simulation and GCM policies}
As introduced in \autoref{subsec:Focus}, the simulation evaluates three temporal resolutions of \ac{gcm} aggregation: \ac{sgcm}, based on the sampled regular room-member group over the study period; \ac{dgcm}, based on the sampled attendees on each day; and \ac{rtgcm}, based on the sampled occupants present at each control timestep.

\begin{figure}[pos=!bt]
    \centering
    \includegraphics[width=0.97\linewidth]{\Pic{png}{OverallFlow}}
    \caption{Overall flow of simulation. Step1 (a)From room user pool, a subgroup of location-log agents is sampled and PCMs are randomly assigned to the sampled location-log agents. Step2(c) Based on the assigned PCMs and location logs,(d) Possibility curve of "No change" in thermal preference id merged based on three types of GCM at each control timestep (d) Setpoints are selected at the point where the aggregated comfort probability is maximized (e) The setpoints are evaluated by returning to the assigned PCMs and extracting the predicted ``No change'' probability from each PCMs}
    \label{fig:ComChara}
\end{figure}

Using the prepared location-log agents and \acp{pcm}, Monte Carlo trials were generated to represent different combinations of observed occupancy compositions and assigned thermal-preference profiles for each target room.
To preserve room-specific usage patterns while varying the simulated occupant mix, each Monte Carlo trial sampled a subgroup of location-log agents from the regular member pool of room $r$ and randomly assigned \acp{pcm} to the selected location-log agents.
The subgroup size $K$ denotes the number of location-log agents sampled from the regular member pool of room $r$ in a simulation condition.
Within a trial, the sampled subgroup size $K$ is fixed for the full simulation period.
However, only a subset of those sampled agents is present at any control timestep, so the realized occupancy count varies according to the observed stay logs and is denoted by $N_{r,t}=|S_{r,t}|$.
The selected location-log agents retained their observed stay logs in their original rooms, while \acp{pcm} were randomly assigned to those selected location-log agents in each trial.
This procedure preserved observed room usage patterns while varying the association between occupancy patterns and comfort-response characteristics.
At each 30-minute control timestep $t$, the subset of sampled location-log agents who are present in room $r$ defined the active set $S_{r,t}$, and each policy selected the corresponding temperature setpoint from its \ac{gcm}.

For each Monte Carlo trial, the three \acp{gcm} described above were formed by averaging the assigned individual ``No change'' probability curves at different temporal resolutions.
For each policy, the selected setpoint was the temperature that maximized the corresponding aggregated ``No change'' probability.
Thus, the three policies share the same setpoint-selection rule and differ only in the update timescale of the represented group: the full study period, each day, or each control timestep.
Let $p_i(T)=P_i(S_{th}=0\mid T)$ denote the ``No change'' probability predicted by the \ac{pcm} assigned to selected location-log agent $i$, where $i$ indexes an individual occupant proxy in the simulation, at \acl{ti} $T$.
For an occupant group $G$, the merged ``No change'' probability curve is defined by averaging the individual curves at each candidate temperature:
\begin{equation}
\bar{p}_{G}(T)
=
\frac{1}{|G|}\sum_{i\in G} p_i(T)
=
\frac{1}{|G|}\sum_{i\in G}
P_i\!\left(S_{th}=0\,\middle|\,T\right).
\label{eq:gcm_probability_curve}
\end{equation}

The group-level setpoint is then selected as
\begin{equation}
T_G^{*}
=
\operatorname*{arg\,max}_{T\in\mathcal{T}_{\mathrm{set}}}
\bar{p}_{G}(T),
\label{eq:gcm_optimal_setpoint}
\end{equation}
where $\mathcal{T}_{\mathrm{set}}$ is the temperature grid of candidate setpoints.
Thus, the common setpoint is selected from the aggregated probability curve, rather than by averaging the individual preferred temperatures.
The policy-specific groups are
\begin{equation}
G_{r}^{\mathrm{Static}} = R_r,
\qquad
G_{r,d}^{\mathrm{Daily}} = A_{r,d},
\qquad
G_{r,t}^{\mathrm{RT}} = S_{r,t},
\label{eq:policy_groups}
\end{equation}
where $R_r$ is the sampled room member group for room $r$, $A_{r,d}$ is the set of sampled location-log agents who appear in room $r$ on day $d$, and $S_{r,t}$ is the set of sampled location-log agents present at timestamp $t$.
The resulting policy setpoints are
\begin{equation}
T_{\mathrm{set}}^{\mathrm{Static}}(t)=T_{G_{r}^{\mathrm{Static}}}^{*},
\qquad
T_{\mathrm{set}}^{\mathrm{Daily}}(t)=T_{G_{r,d(t)}^{\mathrm{Daily}}}^{*},
\qquad
T_{\mathrm{set}}^{\mathrm{RT}}(t)=T_{G_{r,t}^{\mathrm{RT}}}^{*},
\label{eq:policy_setpoints}
\end{equation}
where $d(t)$ is the calendar day containing timestamp $t$.


In addition to these three group-model-based policies, a fixed \qty{24}{\celsius} setpoint is used as an additional baseline.
For each room and subgroup size, \qty{1000} Monte Carlo trials were formed by combining sampled location-log agents and assigned \acp{pcm}, and one-month \ac{occ} simulations based on the three \acp{gcm} were conducted.


\subsection{Evaluation Metrics}
\label{sec:EvalMet}

\subsubsection{Occupancy-related evaluation parameters}
\label{sec:OccuMetrics}
\autoref{table:OccuMetrics} in Appendix~\ref{app:OccupancyStudy} summarizes occupancy-related literature and metrics after COVID-19.
While many studies emphasize occupancy count and related parameters, this study also considers occupant-composition turnover.
Let $\mathcal{T}_{r,d}$ denote the ordered set of analysis timestamps for room $r$ on day $d$, and let $S_{r,t}$ denote the set of sampled location-log agents present in room $r$ at timestamp $t$.
The realized occupancy count is $N_{r,t}=|S_{r,t}|$.

The subgroup size $K$ is the number of sampled location-log agents in a simulation condition, whereas $N_{r,t}$ is the number of those location-log agents actually present at timestamp $t$.

\paragraph{Realized occupancy count}
\begin{equation}
N_{r,t} \;=\; |S_{r,t}|.
\label{eq:Nt}
\end{equation}

\paragraph{Daily mean occupancy}
The daily mean occupancy is
\begin{equation}
\bar{N}_{r,d}
=
\frac{1}{|\mathcal{T}_{r,d}|}
\sum_{t\in\mathcal{T}_{r,d}} N_{r,t}.
\label{eq:mean_occ_daily}
\end{equation}

For comparison across subgroup sizes, the instantaneous occupancy density and daily attendance ratio are defined as
\begin{equation}
\rho_{r,t}^{\mathrm{occ}}=\frac{N_{r,t}}{K},
\qquad
\rho_{r,d}^{\mathrm{att}}=\frac{|A_{r,d}|}{K},
\label{eq:occupancy_and_attendance_ratios}
\end{equation}
where $A_{r,d}$ is the set of sampled location-log agents who attend room $r$ at least once on day $d$.

% \paragraph{Occupancy variability}
% To capture within-day fluctuations in room load, the daily standard deviation of realized occupancy count is used:
% \begin{equation}
% \sigma_{N,r,d}
% =
% \sqrt{
% \frac{1}{|\mathcal{T}_{r,d}|-1}
% \sum_{t\in\mathcal{T}_{r,d}}
% \left(N_{r,t}-\bar{N}_{r,d}\right)^2
% }.
% \label{eq:occ_var}
% \end{equation}

\paragraph{User turnover rate}
The short-horizon replacement of attendees is quantified using a rolling Jaccard turnover, denoted as \ac{utr}.
For timestamp-level analysis, \ac{utr} is computed by comparing the current occupant set with the set observed 30 minutes earlier:
\begin{equation}
\mathrm{UTR}_{r,t}^{30}
=
1
-
\frac{
|S_{r,t}\cap S_{r,t-30}|
}{
|S_{r,t}\cup S_{r,t-30}|
}.
\label{eq:rolling_turnover}
\end{equation}
Here, $S_{r,t-30}$ denotes the occupant set in the same room 30 minutes before timestamp $t$.
When both the current and lagged sets are empty, the timestamp-level turnover is treated as undefined.
Higher $\mathrm{UTR}_{r,t}^{30}$ indicates stronger short-horizon replacement of attendees.
For the time-of-day profiles, timestamp-level \ac{utr} values are summarized by room, weekday, and time of day.

% For daily summaries, timestamp-level \ac{utr} values are averaged within each room-day:
% \begin{equation}
% \mathrm{UTR}_{r,d}
% =
% \frac{1}{|\mathcal{V}_{r,d}|}
% \sum_{t\in\mathcal{V}_{r,d}}
% \mathrm{UTR}_{r,t}^{30},
% \label{eq:utr_daily}
% \end{equation}
% where $\mathcal{V}_{r,d}$ is the set of timestamps on day $d$ for which the 30-minute turnover is defined.
% The daily \ac{utr} captures how much the present group composition changes over short horizons, even when the occupancy count is similar.


\subsubsection{Comfort-related outcomes}
\label{subsec:comfMetrics}

\paragraph{Policy setpoint}
Let $T_{\mathrm{set}}^{\mathrm{policy}}(t)$ denote the setpoint at time $t$ generated by $\mathrm{policy}\in\{\mathrm{Static},\mathrm{Daily},\mathrm{RT}\}$.
Here, $\mathrm{policy}\in\{\mathrm{Static},\mathrm{Daily},\mathrm{RT}\}$ corresponds to the \ac{sgcm}, \ac{dgcm}, and \ac{rtgcm} policies, respectively.

\paragraph{Occupant-level comfort probability}
For each present member $i\in S_{r,t}$, the setpoint generated by a policy is evaluated by the predicted ``No change'' probability $p_i(T_{\mathrm{set}}^{\mathrm{policy}}(t))$ from the assigned \ac{pcm}.

\paragraph{Room-day comfort probability}
The room-day comfort probability is then computed as the person-time average:
\begin{equation}
\mathrm{CP}_{r,d}^{\mathrm{policy}}
=
\frac{
\sum_{t\in\mathcal{T}_{r,d}}
\sum_{i\in S_{r,t}}
p_i\!\left(T_{\mathrm{set}}^{\mathrm{policy}}(t)\right)
}{
\sum_{t\in\mathcal{T}_{r,d}}
|S_{r,t}|
}.
\label{eq:comfort_probability_daily}
\end{equation}
Here, $\mathcal{T}_{r,d}$ is the ordered set of analysis timestamps for room $r$ on day $d$.
Timestamps with no present location-log agents do not contribute person-time comfort scores.

\paragraph{Mean comfort probability}
For comparison across the simulation period, daily room-level comfort probabilities were averaged over the analysis days and Monte Carlo trials for each room, subgroup size, and policy:
\begin{equation}
\overline{\mathrm{CP}}_{r,K}^{\mathrm{policy}}
=
\frac{1}{|\mathcal{M}_{r,K}|}
\sum_{m\in\mathcal{M}_{r,K}}
\frac{1}{|\mathcal{D}_{r,m}|}
\sum_{d\in\mathcal{D}_{r,m}}
\mathrm{CP}_{r,d,m}^{\mathrm{policy}} .
\end{equation}
where $m$ is the Monte Carlo trial index, $\mathcal{M}_{r,K}$ is the set of Monte Carlo trials for room $r$ and subgroup size $K$, $\mathcal{D}_{r,m}$ is the set of valid analysis days for room $r$ in trial $m$, and $\mathrm{policy}\in\{\mathrm{Static},\mathrm{Daily},\mathrm{RT},24\,^\circ\mathrm{C}\}$ denotes the evaluated control policy.

\paragraph{Comfort probability improvement}
The improvement relative to the fixed \qty{24}{\celsius} baseline is computed as
\begin{equation}
\Delta \overline{\mathrm{CP}}_{r,K}^{\mathrm{policy}}
=
\overline{\mathrm{CP}}_{r,K}^{\mathrm{policy}}
-
\overline{\mathrm{CP}}_{r,K}^{24\,^\circ\mathrm{C}} .
\label{eq:comfort_probability_improvement}
\end{equation}

\subsubsection{Control-related outcomes}

\paragraph{Setpoint adjustment magnitude}
To quantify setpoint adjustment requirements under the real-time policy, the setpoint adjustment from the previous control timestep is defined as
\begin{equation}
\delta T_{r,t}^{\mathrm{RT}}
=
T_{\mathrm{set}}^{\mathrm{RT}}(t)
-
T_{\mathrm{set}}^{\mathrm{RT}}(t-\Delta t),
\label{eq:rt_setpoint_step}
\end{equation}
where $\Delta t=\qty{30}{\minute}$ is the control timestep, and $\delta T_{r,t}^{\mathrm{RT}}$ is the setpoint change required by the real-time policy from the previous control timestep.
The adjustment magnitude is evaluated as $|\delta T_{r,t}^{\mathrm{RT}}|$.
Because this analysis focuses on actual setpoint adjustments, the daily mean adjustment magnitude is computed only over timesteps with a non-negligible setpoint change.
\begin{equation}
\mathcal{C}_{r,d}^{+}
=
\left\{
t\in\mathcal{C}_{r,d}
\;:\;
\left|\delta T_{r,t}^{\mathrm{RT}}\right| \geq \varepsilon
\right\},
\qquad
\varepsilon=\qty{0.05}{\celsius}.
\label{eq:changed_step_set}
\end{equation}
\begin{equation}
\overline{|\delta T|}_{r,d,+}^{\mathrm{RT}}
=
\frac{1}{|\mathcal{C}_{r,d}^{+}|}
\sum_{t\in\mathcal{C}_{r,d}^{+}}
\left|\delta T_{r,t}^{\mathrm{RT}}\right|,
\label{eq:rt_mean_abs_step_when_changed_daily}
\end{equation}
where $\mathcal{C}_{r,d}$ is the set of timestamps on day $d$ in room $r$ for which both $T_{\mathrm{set}}^{\mathrm{RT}}(t)$ and $T_{\mathrm{set}}^{\mathrm{RT}}(t-\Delta t)$ are defined, $\mathcal{C}_{r,d}^{+}$ is the subset of valid control timesteps with a non-negligible setpoint change, and $\overline{|\delta T|}_{r,d,+}^{\mathrm{RT}}$ is the mean absolute adjustment magnitude when the setpoint changes.

\paragraph{Action-needed window ratio}
For the Results analysis of control action frequency, the main actionable metric is the 30-minute action-needed window ratio.
For the 30-minute decision-window analysis, the real-time setpoint is first rounded to the nearest \qty{0.5}{\celsius}, denoted by $T_{\mathrm{set},0.5}^{\mathrm{RT}}(t)$.
A 30-minute window is counted as requiring action when at least one rounded real-time setpoint update within the window is \qty{0.5}{\celsius} or larger:
\begin{equation}
\mathrm{AWR}_{r,d}^{30}
=
\frac{1}{|\mathcal{W}_{r,d}^{30}|}
\sum_{w\in\mathcal{W}_{r,d}^{30}}
\mathbb{1}
\left(
\max_{t\in\mathcal{C}_{r,d,w}}
\left|
T_{\mathrm{set},0.5}^{\mathrm{RT}}(t)
-
T_{\mathrm{set},0.5}^{\mathrm{RT}}(t-\Delta t)
\right|
\geq \qty{0.5}{\celsius}
\right),
\label{eq:action_needed_window_ratio}
\end{equation}
where $\mathbb{1}(\cdot)$ is the indicator function, $\mathcal{W}_{r,d}^{30}$ is the set of 30-minute decision windows for room $r$ on day $d$, $\mathcal{C}_{r,d,w}$ is the set of valid within-window consecutive setpoint-update timestamps in window $w$, and $\mathrm{AWR}_{r,d}^{30}$ is the share of 30-minute windows requiring at least one actionable real-time setpoint update.

The defined metrics are used for three comparisons: mean comfort probability and improvement relative to the fixed \qty{24}{\celsius} baseline across subgroup sizes and \ac{gcm} resolutions; real-time setpoint adjustment requirements through mean absolute setpoint adjustment; and the relationship between \ac{utr} and the 30-minute action-needed window ratio.
Representative \ac{utr} transitions by room and weekday are then used to interpret the occupancy-dynamics patterns associated with higher action-needed ratios.

% Removed

% \subsubsection{Personal Comfort Model (PCM)}

% Individual \acp{pcm} are developed for each participant to predict thermal preference based on indoor air temperature.
% Environmental variables are sourced from compact indoor environment sensors (indoor air temperature (\acs{ti}).
% The baseline \ac{pcm} are a multi-class logistic classification model, used in similar agentic-based thermal comfort simulations~\cite{jung_comparative_2019, jung_energy_2020}.

% %study by Daum et al.~\cite{daum_personalized_2011}, widely adopted in prior studies~\cite{jung_comparative_2019,lee_implementation_2019,zhang_impact_2023,meimand_personal_2024,topak_collective_2023,gupta_incentive-based_2018}.

% Specific model inputs and outputs are defined as follows:
% \begin{itemize}
%     \item \textbf{Input vector $\mathbf{x}$:} The primary feature is indoor air temperature (\acs{ti}) obtained from the compact sensors nearest to the occupant's zone.
%     \item \textbf{Output vector $\mathbf{y}$:} The model predicts the probability distribution over three thermal preference classes: $S_{th} \in \{+1 (\text{Warmer}), 0 (\text{No change}), -1 (\text{Cooler})\}$.
% \end{itemize}

% \begin{equation}
% P(S_{th} = s \mid t_i)\,.
% \label{eq:pcm_probability}
% \end{equation}

% \begin{figure}[pos=!bt]
%     \centering
%     \includegraphics[width=0.5\linewidth]{\Pic{png}{ex_comfortmodel_example}}
%     \caption{Individual comfort models from comfort surveys}
%     \label{fig:ComChara}
% \end{figure}

